\documentclass[11pt]{article}
\usepackage[utf8]{inputenc}
\usepackage[french]{babel}
\usepackage[left=2cm,right=2cm,top=2cm,bottom=2cm]{geometry}
\usepackage[style=numeric, sorting=none]{biblatex}

\title{La reconnaissance faciale dans les images}
\author{Lucresse BEBEY DOOH, Adam TRAD, Seif BAHGAT et Cérine BEN AHMED}
\date{1 Décembre 2025}

\begin{document}
\maketitle

\section{Introduction}
Alors qu’il y a quelques années encore, notre société française s’étonnait de la rigueur de l’état de surveillance qu’impose la 
Chine à sa population ; tant soit-il par le biais de détecteur d’infraction, que de possibilité de payer à la simple exposition 
de son visage à une caméra. Un basculement s’est opéré quand les Jeux Olympique de 2024 ont été prétexte à faire voter une loi 
permettant le déploiement de technologies de vidéo-surveillance par algorithme, marquant les premiers pas vers l’analyse de la 
personne humaine par une machine dans le pays.\\
C’est dans un contexte où l’état de la reconnaissance faciale dans les images ne fait que progresser que la recherche lui 
consacre des études visant à expliquer ses différents enjeux : matériels, éthiques, juridiques, etc.
Ce sont tant de questions qui se posent légitimement à la mention de ces sujets ; quels moyens physiques sont nécessaires au 
développement d’une telle technologie, quelles sont ses limites, sur quoi est-elle entrainée, selon quelles conjectures, 
peuvent-elles être neutres, et quel cadre est posé par la loi ?\\
C’est pour essayer d’apporter des réponses à ces questions que nous allons proposer dans cette synthèse quelques articles 
soigneusement choisis et directement relatif à cette problématique. Cette sélection s’est réalisée en suivant les étapes de la 
revue de littérature, d’abord en choisissant nos sources d’informations, elles ont été les suivante : HAL, arXiv, DBLP et Google 
Scholar ; sur ces bases de données nous avons recherché des mots clés tels que « reconnaissance faciale », « reconnaissance 
faciale image », « facial recognition » et « face recognition ».\\
Sur la base de ces recherches nous avons pu fournir notre bibliothèque Zotero d’une quarantaine d’articles en apparence 
cohérents. Pour parvenir à ne retenir que quelques articles nous avons passé cette bibliothèque au crible en regardant dans un 
premier temps si les titres paraissaient pertinents, puis si les résumés l’étaient, jusqu’à pouvoir sélectionner chacun un 
article qui nous intéressait.

\section{Présentation des articles choisis}

\subsection{\textit{Article 1: A Comprehensive Review of Face Recognition Techniques, Trends, and Challenges}, }

Ici les auteurs veulent comprendre l’évolution rapide de la reconnaissance faciale et la hausse  d’usage dans la sécurité, les 
paiements, la santé et la surveillance publiques. Ils  montrent que les systèmes doivent gérer des visages réels avec des 
variations de pose, d’éclairage/ de lumière, d’expression et d’occlusion(par exemple, un visage tourné à 60 degrés ou un visage 
couvert par une écharpe provoque une chute importante du taux de reconnaissance) . On se rend très vite compte que les besoins 
augmentent plus vite que la robustesse des modèles.

La contribution principale est une classification détaillée des méthodes existantes pour l'analyse des images. Les auteurs 
distinguent trois approches :
\begin{itemize}
    \item \textbf{Les méthodes 2D (basées sur l’apparence)} utilisent uniquement les images    classiques, faciles à capturer. 
Elles transforment un visage en vecteurs de caractéristiques, comme PCA ou Eigenfaces. Un exemple typique est la détection de 
visage avec le détecteur Viola Jones, qui fonctionne bien pour les visages frontaux.
    \item \textbf{Les méthodes 3D} s’appuient sur des capteurs qui enregistrent la profondeur, comme les scanners utilisés dans 
les bases BU 3DFE ou Gavab. Elles résistent mieux aux changements de pose, mais sont coûteuses et difficilement accessibles.
    \item \textbf{Les méthodes hybrides } qui mélangent informations globales (comme la forme du visage) et informations locales 
(comme les yeux ou la bouche). Certains travaux combinent Gabor et LDA pour obtenir de meilleures performances en condition de 
lumière difficile.
\end{itemize}
Les auteurs analysent aussi les techniques modernes liées au deep learning comme : 
\begin{itemize}
    \item \textbf{Les CNN} qui sont utilisés pour extraire des caractéristiques faciales complexes
    \item \textbf{Les architectures hybrides} (tel que CNN + SVM, CNN + LSTM, modèles multimodaux)
    \item \textbf{Les fonctions de pertes modernes} comme ArcFace pour améliorer la séparation entre identités.
\end{itemize}
La méthodologie repose sur une revue organisée de la littérature récente. Les auteurs classifient leurs travaux selon trois 
approches : apparence globale, géométrie locale et hybride. Ils évaluent chaque étape du pipeline, de la détection du visage à 
l’extraction des caractéristiques et finalement la classification. Par exemple, HOG permet d’extraire des gradients utiles pour 
repérer les contours du visage. Ils comparent une fois de plus les performances selon les différentes variations de pose, comme 
un visage souriant, un visage dans l’ombre ou un visage partiellement couvert par des lunettes. Ils utilisent des bases de 
données 2D et 3D reconnues comme \textbf{LFW} pour les images naturelles et \textbf{IJB-S} pour les vidéos de surveillance. 
Cette approche permet au lecteur de comprendre les progrès et les limites de chaque technique en conditions réelles.

Cet article met en avant quelques limites importantes encore non résolues actuellement de la reconnaissance faciale :
\begin{itemize}
    \item \textbf{Biais des données d’entraînement} : Les bases de données présentent souvent un manque de diversité            
ethnique. Cela crée des défaillances pour certaines populations et ce déséquilibre réduit la fiabilité des          systèmes.
    \item \textbf{Variations extrêmes de pose} : Certains modèles échouent encore lorsque les visages de personnes sont très 
inclinés voire tournés dans un angle peu favorable. Cela est dû au passage d’un visage nonfrontal à un visage frontal qui reste 
encore un problème majeur.
    \item \textbf{Illumination, expressions et occlusions} : Malgré les modèles récents, nous avons toujours des pertes 
d’informations et des erreurs de classification constantes lorsque nous avons plusieurs facteurs combinés, tel que le port de 
lunettes, ou bien faire une grimace.
    \item \textbf{Manque de données 3D accessibles} : Étant donné que les bases 3D sont rares et coûteuses à produire, ce manque 
limite l’évolution des approches 3D.
    \item \textbf{Défis dans la reconnaissance vidéo} : Les auteurs montrent aussi que la reconnaissance vidéo reste difficile. 
Une vidéo de surveillance contient souvent des visages flous, des résolutions basses ou des séquences où la même personne 
n’apparaît que quelques frames. Par exemple, un visage vu de loin dans une caméra de rue donne moins d’information qu’une image 
frontale en haute résolution.
\end{itemize}

\subsection{\textit{Article 2: La Reconnaissance Faciale : Entre Promesse et Péril}, Thierry Ménissier}

Les dispositifs de reconnaissance faciale (DRF) se répandent à grande vitesse dans nos sociétés. Ils promettent plus 
d’efficacité, plus de sécurité, et une gestion simplifiée de nombreux services. Mais derrière ces promesses se cachent aussi des 
questions très sérieuses sur nos libertés, notre vie privée et, plus largement, sur la manière dont nous voulons vivre en 
démocratie.

Les DRF ne se contentent pas d’identifier les individus sans contact ni accord explicite. Ils vont parfois beaucoup plus loin : 
certaines technologies cherchent à analyser nos comportements, nos émotions ou même nos opinions. Et une fois croisées avec 
d’autres données personnelles bancaires, médicales, administratives, elles peuvent produire des profils extrêmement détaillés de 
chacun d’entre nous, parfois à notre insu.

Tout cela ouvre la porte à un modèle de société où la surveillance devient presque invisible, mais omniprésente : un véritable 
“panoptique numérique” dans lequel l’anonymat, pourtant essentiel à la liberté d’expression et à la vie démocratique, s’efface 
peu à peu. Dans les “villes intelligentes”, le risque est que le contrôle devienne automatique, continu, et quasiment impossible 
à contourner.

Sur le plan éthique, les DRF sont souvent évalués selon une logique utilitariste : on met les bénéfices d’un côté, les risques 
de l’autre, et on calcule. Mais cette vision est trop limitée. L’auteur défend une approche plus large qui prend aussi en compte 
des perspectives déontologiques (le respect des droits), arétaïques (les vertus, le type de société que l’on veut construire) et 
axiologiques (les valeurs fondamentales).

Face à ces enjeux, deux moyens d’action apparaissent indispensables : d’une part, garantir une information claire et un 
véritable consentement avant toute utilisation de la reconnaissance faciale ; d’autre part, apprendre aux citoyens à refuser les 
identifications biométriques quand elles ne sont ni utiles ni justifiées.

Au fond, la reconnaissance faciale pose une question simple mais essentielle : sommes-nous prêts à troquer un peu de confort et 
de sécurité contre une part de notre liberté ? Ou préférons-nous préserver l’anonymat, l’autonomie et la dignité qui font de 
nous des individus libres ?

\subsection{\textit{Article 3: Facial recognition technology and protection of fundamental rights}, Raphaël Déchaux}
Raphaël Déchaux examine la technologie de reconnaissance faciale (FRT) et son impact sur la sauvegarde des droits fondamentaux 
en Europe. Le sujet se structure autour de trois axes majeurs : l'explication et le mécanisme de la reconnaissance faciale, les 
menaces qu'elle fait peser sur les droits de l'homme, et les interventions réglementaires européennes destinées à réguler son 
utilisation.

Dans une première section, l'écrivain souligne que les images faciales servent d'identifiant biométrique distinctif, utilisé par 
des systèmes d'identification automatique. Il identifie différentes modalités d'usage : la vérification (un visage confronté à 
une identité déclarée), l'identification (un visage confronté à une base de données), ainsi que la reconnaissance faciale 
automatisée en temps réel. On note également l'émergence de nouvelles applications telles que la classification (âge, sexe, 
origine), l'analyse des émotions et les systèmes dits « physiognomoniques », qui prétendent inférer un comportement ou une 
intention à partir des caractéristiques du visage, des pratiques fréquemment qualifiées de pseudo-scientifiques. Aujourd'hui, la 
technologie est largement mise en œuvre tant dans le secteur public que privé.

La deuxième partie présente les risques. D’abord, Déchaux souligne les dysfonctionnements possibles, faux positifs et faux 
négatifs, souvent liés à la qualité des images ou à des données d’entraînement incomplètes dont les conséquences sont lourdes, 
notamment en matière policière. Ensuite, il met en évidence cette évolution d’un usage administratif de la FRT vers une logique 
de surveillance, où plusieurs décisions des juridictions européennes alertent sur les dérives possibles. S’ajoutent à cela des 
enjeux majeurs de discrimination (certaines catégories de population moins bien reconnues), de protection de la vie privée et 
d’usage illégal de données biométriques, comme le montre le scandale Clearview AI. Enfin, selon l’auteur, la FRT peut également 
menacer des libertés fondamentales comme la liberté d’expression, de réunion ou d’association, notamment si elle est utilisée 
pour surveiller des populations ou des opposants politiques.

Enfin, la troisième partie passe en revue la réponse réglementaire européenne. En l’absence d’un cadre national organisé dans de 
nombreux États membres, l’auteur montre l’intérêt d’une régulation au niveau international et européen. Le Conseil de l’Europe 
défend des principes dont la transparence, le droit de comprendre les décisions automatisées ou le droit de ne pas être soumis à 
de telles décisions. L’Union européenne, de son côté, moultne progressivement un cadre autour de ces technologies à travers 
plusieurs textes mais surtout l’AI Act. Celui-ci classe les systèmes de reconnaissance faciale selon différentes catégories de 
risques : certains usages sont jugés « à haut risque » et strictement encadrés (identification biométrique à distance, 
catégorisation sensible, reconnaissance émotionnelle), tandis que d’autres sont purement prohibés, à l’instar du recours à la 
création de bases d’images faciales par scraping massif ou de l’identification biométrique en temps réel dans l’espace public.

En résumé, Déchaux souligne la demande d’un encadrement juridique étroit pour préserver les droits fondamentaux face à un 
progrès technologique qui atteint une grande vélocité. La reconnaissance faciale peut éventuellement rendre service dans 
certaines situations mais elle emporte de lourds risques pour la vie privée, la non-discrimination et les libertés publiques. 
C’est la raison pour laquelle l’Europe tente d’arbitrer entre innovation, sécurité et droits humains.

\subsection{\textit{Article 4: Des visages invisibles. Vers une histoire des images de la reconnaissance faciale de Roland 
Meyer}, traduit par Florence Rougerie}
Cette article conte le fil de l’histoire de la reconnaissance faciale dans les images, de son origine à l’état dont on pouvait 
la trouver en 2020, date de parution de ce texte. Il revient chronologiquement sur les grandes étapes de l’évolution de ce 
domaine, en orientant le lecteur vers une problématique annonçant déjà le ton critique : la reconnaissance faciale telle que « 
pratiquée » par les machines aura beau s’efforcer d’y parvenir, elle ne reflétera jamais la complexité de la perception d’un 
visage par une personne.
Cette chronologie débute par la citation « Rien n’est moins personnel que le visage », extraite des carnets de Deleuze et Parnet 
(1977) ; et est tout du long mise face à la théorie de la « visagéité » de Deleuze et Guattari, laquelle théorie présente le 
visage par le prisme de sa binarité, un trou noir symbolisant les éléments distinctifs, et un mur blanc renvoyant aux 
interprétations que l'on en fait ; ces deux là formant une machine sociale donnant au visage ses codes.

La première manifestation de « reconnaissance faciale » se retrouve dans le fait de placarder des avis de recherche d’une 
personne. Les prémisses de ce domaine se fondent sur la volonté de la police de pouvoir afficher et identifier des individus 
plus facilement, ici donc, la circulation d’image « décide du destin d’êtres humains ».
Il est à ce moment encore difficile d’aboutir à un système de reconnaissance par l’image fonctionnel car les moyens techniques 
sont peu avancés et la transposition de cette complexité humaine parait alors être une tâche herculéenne.

Il faudra attendre le projet Computer Physiognomy mené par la Nippon Electric Company et exposé au cours de l’exposition 
universelle d’Osaka de 1970 pour que de premières bases soient posées. Des visiteurs sont invités à se faire photographier puis 
fournir ce cliché à un système de traitement qui déterminera ensuite à quelle célébrité parmi une liste de 7 ils ressemblent le 
plus.
Sur la base de ces expériences Takeo Kanade, ingénieur japonais met au point l’un des premiers systèmes fonctionnels de 
reconnaissance faciale, et ce en pixelisant les images et codifiant les écart que l’on trouve entre les différents éléments du 
visage, la machine sait alors en principe où trouver une telle partie.

Meyer souligne que la science de l’analyse des visages humains résultant à ses origines d’une volonté coloniale, nul algorithme 
de reconnaissance faciale ne peut exister sans être biaisé.
Dans les années 70 à 80, différentes sciences connaissent une nouvelle façon d’aborder leurs études avec le « tournant cognitif 
», modélisant la pensée humaine comme un traitement de données. 
Il faudra attendre que soit passé l’hiver de l’IA des années 80 et qu’arrivent les années 90 pour qu’un nouveau tournant voit le 
jour ; la reconnaissance faciale n’est alors plus tant envisagée comme une façon de reproduire la perception humaine d’un visage 
mais plutôt par l’exploitation des informations figurant sur les images. En 1991, Matthew Turk et Alex Pentland du MIT, dédient 
un essai à cette thèse : « Eigenfaces for Recognition ».
Dès lors la mise en place d’un sytème de reconnaissance faciale nécessite une base de données de visages alternatifs et 
standardisés (eigenfaces), desquels le premier enseignement et d’apprendre à reconnaitre un visage. Les visages ne sont alors 
plus composés de caractéristiques mais de « structures de données ».

Les évènements du 11 septembre marquent un tournant dans la course au développement d’une technologie de reconnaissance faciale 
de par la nécessité de généraliser l’équipement biométrique pour faciliter le contrôle aux frontières.
Enfin, dès les années 2010, la mise en relation de deux technologies développées séparément : les réseaux neuronaux artificiels 
et les réseaux sociaux va permettre d’amorcer la notion de reconnaissance faciale telle qu’on la connait.
Meyer finit par conclure son article en rappelant au lecteur qu’en participant à l’activité d’un réseau social à partir de post, 
commentaire, ou encore identification ; chacun nourri ces bases de données d’informations personnelles, soient-elles les nôtres 
ou celles de nos proches.


\section{Analyse comparative}
Les cinq résumés sont abordés d’une façon complémentaire, permettant une comparaison riche entre approches techniques, 
juridiques et philosophiques. Certes, ils divergent dans certains éléments, mais tous convergent vers une idée centrale : la 
reconnaissance faciale est une technologie à fort potentiel, mais dont les risques pourraient compromettre définitivement les 
droits et les libertés substantiels.

\subsection{Similarités}
Chacun des articles, Déchaux et Ménissier d’une part, et l’article éthique et les deux articles techniques d’autre part, précise 
que la reconnaissance faciale connaît un marché en forte expansion et se diversifie dans différents champs d’usage : sécurité, 
surveillance, paiements, santé ou smart cities. Tous s’accordent également sur l’existence de risques conséquents. Tandis que 
les articles de Déchaux et Ménissier concernent la liberté fondamentale (vie privée, anonymat ou liberté d’expression), ils sont 
les seuls à mettre en avant les risques liés à l’imperfection des systèmes (reconnaissance erronée, biais des datasets, effets 
de la lumière, des poses, des occlusions), alors que dans les articles techniques, ce sont les limites et manques de performance 
qui apparaissent comme des dangers propres au fonctionnement de la technologie. Ces limites techniques s’accordent avec les 
traditions éthiques et politiques, car elles peuvent mener à des injustices, notamment dans le traitement policier. Enfin, tous 
reconnaissent le besoin d’encadrement, qu’il soit légal ou de nature différente.\\
Les grandes différences apparaissent dans la nature même des analyses. Les deux articles techniques ont une visée descriptive et 
scientifique : ils typifient les méthodes (2D, 3D, hybrides), décrivent les architectures (CNN, modèles multimodaux) et mesurent 
les performances sur des bases reconnues. Ils ne cherchent pas à juger la technologie, mais à la situer dans ses progrès et ses 
limites, expliquant véritablement comment fonctionne la reconnaissance faciale. À l’inverse, Ménissier et l’article éthique 
proposent une lecture normative. Ménissier mobilise des concepts philosophiques comme le panoptique, le biopouvoir, et 
l’identité/personnalité numérique. Il montre que la reconnaissance faciale ne se limite pas à reconnaître : elle personnifie en 
recoupant le visage avec des big data (fiscales, médicales, sociales), ce qui en fait un danger pour la démocratie et 
l’anonymat. L’article éthique, lui, critique la vision purement utilitariste et demande une éthique plus large, prenant en 
compte les droits, les valeurs et les impacts sociaux.


\subsection{Différences / Controverses}
Les grandes différences apparaissent dans la nature même des analyses. Les deux articles techniques ont une visée descriptive et 
scientifique : ils typifient les méthodes (2D, 3D, hybrides), décrivent les architectures (CNN, modèles multimodaux) et mesurent 
les performances sur des bases reconnues. Ils ne cherchent pas à juger la technologie, mais à la situer dans ses progrès et ses 
limites, expliquant véritablement comment fonctionne la reconnaissance faciale. À l’inverse, Ménissier et l’article éthique 
proposent une lecture normative. Ménissier mobilise des concepts philosophiques comme le panoptique, le biopouvoir, et 
l’identité/personnalité numérique. Il montre que la reconnaissance faciale ne se limite pas à reconnaître : elle personnifie en 
recoupant le visage avec des big data (fiscales, médicales, sociales), ce qui en fait un danger pour la démocratie et 
l’anonymat. L’article éthique, lui, critique la vision purement utilitariste et demande une éthique plus large, prenant en 
compte les droits, les valeurs et les impacts sociaux.\\
Les articles techniques ont pour objet principal la fiabilité et la robustesse : diversité des données, limites des bases 3D, 
difficultés de la reconnaissance vidéo. Ils montrent que la technologie pose avant tout des défis scientifiques. À l’inverse, 
Ménissier pense que le véritable problème n’est pas seulement la défaillance technique mais le changement qu’impose cette 
technologie : fin de l’anonymat, surveillance secrète, confusion entre identités civile et numérique. Déchaux participe de ce 
raisonnement en reliant ces risques aux enjeux démocratiques comme la légitimité policière ou les libertés publiques et aux 
décisions judiciaires européennes.


\subsection{Un tout cohérent}
Si chaque résumé donne une vision différente du problème, en fin de compte, les cinq résumés forment un tout cohérent. Les 
articles techniques expliquent pourquoi la reconnaissance faciale peut être inefficace. Déchaux montre, à partir de ces 
impuissances, les enjeux juridiques et politiques que cela représente. Ménissier et l’article éthique explicitent comment la 
mise en œuvre de cette technologie modifie notre rapport à la liberté, à l’espace public et à l’identité. Si pour les premiers 
la technologie est un succès ou un échec, pour les autres elle pose des questions éthiques fondamentales.

\section{Conclusion}

Cette synthèse a permis de confirmer que la \textbf{reconnaissance faciale dans les images} est un sujet d'une richesse et d'une 
complexité remarquables qui mettent en avant les avancées techniques et des enjeux sociétaux fondamentaux.

Les progrès de l'Intelligence Artificielle ont permis à la reconnaissance faciale de franchir des étapes décisives. Comme l'a 
montré l'Article 1, l'adoption des \textbf{réseaux neuronaux du CNN} et des techniques d'optimisation comme la 
\textbf{modélisation 3D} a considérablement augmenter la robustesse et la précision des systèmes. L'évolution de ce domaine, 
retracée par Meyer (Article 4), prouve que la quête d'une identification par l'image initiée il y a des décennies a aujourd'hui 
atteint un niveau de technique remarquable.

En conclusion, la RF se positionne comme un excellent exemple d'innovation technologique pour lequel le \textbf{cadre éthique et 
légal} doit évoluer rapidement pour rattraper le progrès technique. La réponse apportée, notamment par l'Europe avec 
l'\textbf{AI Act} (Article 3), vise à établir un équilibre délicat entre l'impératif de sécurité et la préservation de la vie 
privée. L'avenir de la reconnaissance faciale ne dépend donc plus uniquement de la résolution des défis techniques comme les 
biais des données d'entraînement mais de la capacité de la société à \textbf{maintenir le contrôle} sur son usage et à 
privilégier la \textbf{liberté} sur la seule commodité.

\end{document}
